\documentclass[11pt]{article}
\usepackage{amsmath,amssymb}
\title{An original two-link fiber: gauge cancellation and a transverse signed fluctuation trace}\author{}\date{}
\begin{document}\maketitle
\paragraph{Original fiber and raw normalization.}
Take two consecutive selected spatial links $u,v\in SU(2)$.
The first link has four fixed spatial staples and its preceding temporal
neighbor equal to $A=e^{i\epsilon\sigma_3}$; the second has four spatial
staples and its following temporal neighbor equal to
$C=e^{-i\epsilon\sigma_3}$. The temporal bath links are the identity.
These fixed staples can be realized by the private links in the adjacent
plaquettes; all other fixed plaquettes remain part of the exterior record.
With $s(U)=1-\frac12\operatorname{ReTr}U$, the free-dependent action is
\[
 S=B[\kappa s(uA^{-1})+\kappa s(vC^{-1})+s(uv^{-1})],\qquad \kappa=5.
\]
Its raw integral contains the factor $e^{-11B}$, not a normalized temporal
kernel. Any fixed-only plaquette contribution $S_{\rm ext}(\epsilon)$ gives
$Z_{\rm full}=e^{-S_{\rm ext}}Z$ and
$V_{\rm full}=S_{\rm ext}+V$. It cancels only in the explicitly
classical-relative fluctuation $F=-\log Z_{\rm full}-V_{\rm full}
=-\log Z-V$. No live exterior action is discarded from $Z_{\rm full}$.

For sufficiently small $\epsilon$, compact pinning and the positive flat
Hessian give a unique global action minimum
\[
 u_*=e^{i\theta\sigma_3},\quad v_*=e^{-i\theta\sigma_3},\qquad
 \kappa\sin(\theta-\epsilon)+\sin2\theta=0,\quad
 \theta'(0)=\frac{\kappa}{\kappa+2}.
\]
Indeed uniqueness at zero follows from the nonnegative pinned action;
compactness excludes other global minimizers away from its neighborhood
for small parameters, and local strict convexity gives uniqueness there.
Symmetry and stationarity identify the displayed minimum.

\paragraph{Complete leading Gaussian around the actual minimum.}
Use right fluctuations
$u=u_*\exp(i x\cdot\sigma/\sqrt B)$,
$v=v_*\exp(i y\cdot\sigma/\sqrt B)$. The exact canonical action Hessian is
\[
 H=\begin{pmatrix}aI&D\\D^T&aI\end{pmatrix},\quad
 a=\kappa\cos(\theta-\epsilon)+\cos2\theta,\quad
 D=-cI-sJ_3,\quad c=\cos2\theta,\ s=\sin2\theta,
\]
where $J_3z=e_3\times z$ (reversing the orientation reverses the immaterial
sign of $J_3$). Its eigenvalues are $a\pm1$, each twice, and $a\pm c$,
each once. Hence
\[
 \det H=(a^2-1)^2(a^2-c^2),\qquad
 Z_{\rm quad}=e^{-V}\frac{2}{\pi B^3}(\det H)^{-1/2}.
\]
The prefactor uses normalized Haar measure with local density
$1/(2\pi^2)$ on each $SU(2)$ factor. Fixed exterior factors multiply this
expression exactly as above. This is a leading quadratic reference,
not an equality to the nonlinear compact integral.

\paragraph{Actual gauge path and the signed cancellation.}
Conjugate every exterior link and both integration variables by
$g(t)=e^{it\sigma_1/2}$. This is an actual site-gauge path (the same
transformation at every site), retaining all holonomy eigenvalues.
Both raw $Z_{\rm full}$ and $V_{\rm full}$ are exactly unchanged.
In the minimum-only translated frame the quadratic Hessian rotates:
\[
 H_t=O_t H O_t^T,\quad O_t=\operatorname{diag}(\operatorname{Ad}_{g(t)},
 \operatorname{Ad}_{g(t)}),\quad
 \dot H=[J,H],\quad \ddot H=[J,[J,H]] .
\]
The harmless sign convention for the adjoint generator does not affect
the following identities. For nonzero small $\epsilon$,
\[
 I=\tfrac12\operatorname{Tr}(H^{-1}\dot H H^{-1}\dot H)>0,\qquad
 \|\dot H\|_{\rm F}^2=4s^2,\qquad I\ge\frac{2s^2}{(a+1)^2}.
\]
Nevertheless
\[
 F_{\rm quad}''=
 \tfrac12\operatorname{Tr}(H^{-1}\ddot H)-I=0
\]
because the determinant is invariant. This retains the acceleration
partner; replacing the curved gauge path by its tangent misses it.
In the fully gauge-aligned fluctuation frame the pulled action and Haar
measure are unchanged, so the corrected score is exactly zero,
even nonlinearly. Thus a positive Fisher term in a minimum-only frame
cannot alone be bounded by the physical change of curvature of a gauge
orbit. The signed trace, or a correctly aligned physical frame, is needed.

\paragraph{One physical, localized transverse change.}
Now vary $\epsilon$, rather than conjugating the records.
The relative center $AC^{-1}=e^{2i\epsilon\sigma_3}$ changes eigenvalues;
this direction is not a gauge orbit. At $\epsilon=0$ direct differentiation
gives
\[
 a=\kappa+1,\quad
 a''=-\frac{4\kappa(\kappa+1)}{(\kappa+2)^2},\quad
 c''=-\frac{4\kappa^2}{(\kappa+2)^2},
\]
\[
 F_{\rm quad,\epsilon\epsilon}
 =-\frac{4(3\kappa^2+5\kappa+3)}{(\kappa+2)^3},\qquad
 V_{\epsilon\epsilon}=\frac{4B\kappa}{\kappa+2}.
\]
For canonical $\epsilon=\zeta/\sqrt B$ and the original $\kappa=5$,
\[
 F_{\rm quad,\zeta\zeta}=-\frac{412}{343B},\qquad
 V_{\zeta\zeta}=\frac{20}{7},\qquad
 |F_{\rm quad,\zeta\zeta}|=\frac{103}{245B}V_{\zeta\zeta}.
\]
This explicitly tests the signed determinant response against the actual
finite-fiber classical Schur energy, rather than a coordinate counting
bound. Fixed-only exterior action cancels from $F$; at the flat record
its Hessian is nonnegative, so including it can only increase the
classical curvature used on the right.

\paragraph{Limits and next estimate.}
The gauge invariance statement is exact for the original compact model.
The numerical signed transverse coefficient above concerns the complete
leading Gaussian only. It is not a proof of a uniform nonlinear remainder,
a growing-block curl estimate, or S94. The reusable next target is the
signed Hessian determinant response on the full relative incidence
complex, with its gauge acceleration and measure partners, bounded by
a physical boundary Schur form. Bounds on Fisher alone do not establish
that statement. Global torus modes have not been integrated out here.
\end{document}
